% answers_short.tex —— 2014 年·广东卷·文科 · 简短答案速查
% 由 tools/make_short_answers.py 自 answers.tex 的 \gksolhead/\gkanswer 提取，
% 并与 2014/2014_数学真题_广东文_答案要点.md 逐题比对后生成，勿手改（改源 MD 后重跑）。

\begin{gknotice}
\textbf{使用说明}：本节只给最终答案便于快速核对，完整推导见解析版。答案为编写者独立推导，\textbf{非官方标准答案}。
\end{gknotice}

\gkpart{选择题}
\begin{center}\begin{tabular}{ccccccccccc}
\toprule
题号 & 1 & 2 & 3 & 4 & 5 & 6 & 7 & 8 & 9 & 10 \\
\midrule
答案 & B & D & B & C & A & C & A & D & D & B \\
\bottomrule
\end{tabular}\end{center}

\gkpart{填空题}
\begin{itemize}[leftmargin=2.6em]
  \item[\textbf{11.}] $5x+y+2=0$。$y'=-5\mathrm e^x$，在 $x=0$ 处斜率为 $-5$，切线为 $y+2=-5x$。
  \item[\textbf{12.}] $\dfrac25$。$\dfrac{\mathrm C_4^1}{\mathrm C_5^2}=\dfrac4{10}=\dfrac25$。
  \item[\textbf{13.}] $5$。由 $a_1a_5=a_3^2=4$ 且各项为正得 $a_3=2$，故 $a_1a_2a_3a_4a_5=a_3^5=32$，所求为 $\log_232=5$。
  \item[\textbf{14.}] $(1,2)$。$\rho\cos\theta=1$ 即 $x=1$；$2\rho\cos^2\theta=\sin\theta$ 两边乘 $\rho$ 得 $2x^2=y$，代入 $x=1$ 得 $y=2$。
  \item[\textbf{15.}] $3$。$AE\parallel DC$ 得 $\triangle AEF\sim\triangle CDF$，相似比 $\dfrac{AE}{DC}=\dfrac{AE}{AB}=\dfrac13$，周长比为其倒数 $3$。
\end{itemize}

\gkpart{解答题}
\begin{itemize}[leftmargin=2.6em]
  \item[\textbf{16.}] (1) $3$。$\sin\left(\dfrac{5\pi}{12}+\dfrac\pi3\right)=\sin\dfrac{3\pi}4=\dfrac{\sqrt2}{2}$，故 $A\cdot\dfrac{\sqrt2}{2}=\dfrac{3\sqrt2}{2}$；(2) $\sqrt6$。$f(\theta)-f(-\theta)=3\sin\theta=\sqrt3$ 得 $\sin\theta=\dfrac{\sqrt3}{3}$，$\cos\theta=\dfrac{\sqrt6}{3}$，所求为 $3\sin\left(\dfrac\pi2-\theta\right)=3\cos\theta=\sqrt6$
  \item[\textbf{17.}] (1) 众数为 $30$，极差为 $21$。$30$ 岁出现 $5$ 次最多；最大 $40$、最小 $19$；(2) 茎 $1$ 叶 $9$；茎 $2$ 叶 $8,8,8,9,9,9$；茎 $3$ 叶 $0,0,0,0,0,1,1,1,1,2,2,2$；茎 $4$ 叶 $0$；(3) $12.6$。平均年龄 $\dfrac{600}{20}=30$，偏差平方和 $252$，方差 $\dfrac{252}{20}=12.6$
  \item[\textbf{18.}] (1) 见解析版；(2) $\dfrac{\sqrt2}{16}$。由 $CF\perp$ 平面 $MDF$ 得 $CF\perp DF$，解得 $PF=\dfrac32$、$CF=\dfrac12$、$DF=\dfrac{\sqrt3}{2}$、$MD=\dfrac{\sqrt6}{2}$；$S_{\triangle CDE}=\dfrac{\sqrt3}{8}$，$V=\dfrac13\cdot\dfrac{\sqrt3}{8}\cdot\dfrac{\sqrt6}{2}=\dfrac{\sqrt2}{16}$
  \item[\textbf{19.}] (1) $2$。令 $n=1$ 得 $a_1^2+a_1-6=0$，取正根 $a_1=2$；(2) $a_n=2n$。方程分解为 $(S_n-n^2-n)(S_n+3)=0$，故 $S_n=n^2+n$，相减得 $a_n=2n$；(3) 见解析版
  \item[\textbf{20.}] (1) $\dfrac{x^2}{9}+\dfrac{y^2}{4}=1$。$c=\sqrt5$，由 $\dfrac ca=\dfrac{\sqrt5}{3}$ 得 $a=3$，$b^2=9-5=4$；(2) $x^2+y^2=13$。设过 $P$ 的切线 $y=kx+m$，相切条件化为 $(9-x_0^2)k^2+2x_0y_0k+(4-y_0^2)=0$，由 $k_1k_2=-1$ 得 $x_0^2+y_0^2=13$（竖直切线情形同样满足）
  \item[\textbf{21.}] (1) 当 $a\geqslant1$ 时，单调递增区间为 $(-\infty,+\infty)$；当 $a<1$ 时，增区间为 $(-\infty,-1-\sqrt{1-a})$ 与 $(-1+\sqrt{1-a},+\infty)$，减区间为 $(-1-\sqrt{1-a},-1+\sqrt{1-a})$；(2) 当 $-\dfrac{25}{12}<a<-\dfrac{7}{12}$ 且 $a\neq-\dfrac54$ 时存在；当 $a\leqslant-\dfrac{25}{12}$、$a=-\dfrac54$ 或 $-\dfrac{7}{12}\leqslant a<0$ 时不存在。依据是 $f(x)-f\left(\dfrac12\right)=\left(x-\dfrac12\right)\varphi(x)$，$\varphi(x)=\dfrac13x^2+\dfrac76x+\dfrac7{12}+a$ 在 $(0,1)$ 上递增，须 $\varphi$ 在 $(0,1)$ 内有零点且零点不为 $\dfrac12$
\end{itemize}
